% Source: https://artofproblemsolving.com/community/c3h1575632_help_needed_please_help_me
% Stable identity: aops-topic-1575632-q01
Let $x$ grams come from the $40\%$ bottle; then $20-x$ grams come from the $70\%$ bottle. The amount of pure acid must be $52\%$ of $20$ grams, that is $10.4$ grams:
\[
0.40x+0.70(20-x)=10.4\quad\Longrightarrow\quad 14-0.3x=10.4\quad\Longrightarrow\quad x=12 .
\]
So she uses $\boxed{12\text{ g of the }40\%\text{ solution and }8\text{ g of the }70\%\text{ solution}}$.

Check: $0.4\cdot12+0.7\cdot8=4.8+5.6=10.4$ grams of acid in $20$ grams. (Alternatively, $52$ is $12$ above $40$ and $18$ below $70$, so the amounts are in the ratio $18:12=3:2$.)
